
\section{An Informal Survery of First-Order Logic}

Recall from our discussion that TFL falls short of accurately depicting some
aspects of arguments and inferences in ordinary English. We repeat our ealier
summary for convenience.

\begin{observation}[Shortcomings of TFL]
  Truth-functional logic has some shortcomings that make it ill-suited to deal
  with some phenomena from classical logic and ordinary English:\\
  \begin{enumerate}
  \item TFL does not support the types of inferences depicted in Aristotle's
    syllogisms.
  \item TFL insists on \emph{excluding the middle} -- every TFL is either true
    or false -- but some adjectives do not quite work like that.
  \item TFL comes with a number of \emph{paradoxes} that make the material
    conditional unsatisfying for symbolizing English sentences: one the one
    hand, negating the whole conditional often gives the wrong conclusions
    about its antecedent sentence; on the other hand, negating the consequent
    sentence opens up the possibility of contradiction when the antecedant
    sentence is false.
  \end{enumerate}
\end{observation}

We will need to wait until future logic courses to deal with (2) and (3)--
they are matters for \emph{non-classical logic} and \emph{modal logic},
respectively--but we can start trying to formalize a system to capture
Aristotle's syllogism results right away. This system is \emph{first-order
logic}, also commonly refered to as \emph{predicate logic}.

\subsection{Artistotle and the Ancient Origins of Predicate Logic}

Let's start with a review of Aristotle's syllogisms. The system is based on
the organization of four \emph{quantifiers} (although that term does not
come until later) into a \emph{square of opposition}.

\begin{definition}[Aristotle's \emph{square of opposition}]
  Accounting for the the effects negation on \emph{universal} (``every'') and
  \emph{existential} (``some'') statements, there are four distinct
  \emph{quantifiers}. We can lay them out in opposition to their \emph{duals}
  (i.e. negations) in the following square. Each distinct quantifier is labeled
  with a vowel (A,E,I,O) so that we can refer to it easily. \\
\twoitems{
  \begin{argument}
    \item[A:] Every A is B.
    \item[ ] 
    \item[ ] 
    \item[O:] Some A is not B.
  \end{argument}
}{
  \begin{argument}
    \item[E:] No A is B.
    \item[ ] 
    \item[ ] 
    \item[I:] Some A is B.
  \end{argument}  
}
\end{definition}

Armed with notations for these four types of sentences, we can lay out the
syllogisms (later named by medeival scholars) into three classes, based on
how they are connected by the so-called \emph{middle term} (B), which can
occur in either the \emph{subject} or \emph{predicate} position of each of the
two premise sentences.

\begin{table}[h]
  \centering
  \label{tab:aristotles-syllogisms}
  \begin{tabular}{|c c|c c c|}
    \hline
    \makecell{Medeival\\Name}&\makecell{Aristotle\\Name}& & &\\
    \hline
    Barbara  &(AAA)& All B are C.     & All A are B.    & So all A are C.     \\
    Celarent &(EAE)& No B are C.      & All A are B.    & So no A are C.      \\
    Darii    &(AII)& All B are C.     & Some A is B.    & So some A are C.    \\
    Ferio    &(EIO)& No B are C.      & Some A is B.    & So some A is not C. \\
    \hline
    Cesare   &(EAE)& No C are B.      & All A are B.    & So no A are C.      \\
    Camestres&(AEE)& All C are B.     & No A are B.     & So no A are C.      \\
    Festino  &(EIO)& No C are B.      & Some A is B.    & So some A is not C. \\
    Baroko   &(AOO)& All C are B.     & Some A is not B & So some A is not C. \\
    \hline
    Datisi   &(AII)& All B are C.     & Some B is A.    & So some A is C.     \\
    Disamis  &(IAI)& Some B is C.     & All B are A.    & So some A is C      \\
    Ferison  &(EIO)& No B are C.      & Some B is A.    & So some A is not C. \\
    Bokardo  &(OAO)& Some B is not C. & All B are A.    & So some A is not C. \\
    \hline
    Camenes  &(AEO)& All C are B.     & No B are A.    & So no A are C.       \\
    Dimaris  &(IAI)& Some C is B.     & All B are A.   & So some A is C.      \\
    Fresison &(EIO)& No C is B.       & Some B is A.   & Some A is not C.     \\
    \hline
  \end{tabular}
  \caption{Artistotle's syllogisms with medeival names}
\end{table}

\subsection{Predicates and Singular Terms}

Recall that the problem with treating these examples in TFL is that TFL does not
have a facility for making distinctions between atomic sentences, aside from
their truth. More precisely, we cannot capture the relationship between ``all
men'' and ``Socrates'' that allows us to make the following inference, analogous
to \emph{Barbara}.

\begin{argument}
\item[ ] All men are mortal. 
\item[ ] Socrates is a man.
\item[\therefore] Socrates is mortal.
\end{argument}

Nor can we capture the fact that ``is/are mortal'' and ``is a man'' are
something like properties that we can apply to things like ``all men'' and
``Socrates''. With first-order logic, we get some rudimentary tools to depict
these phenomena formally.

Let's start with some a informal review of the phenomena.

Consider first some simple sentences that differ as little as possible.

\begin{argument}
\item[ ] John is reading.
\item[ ] Mary is reading.
\item[ ] Fred is reading.
\item[ ] Susie is reading.
\end{argument}

We can see a pattern here involving two components.

First, there is whatever the expressions ``John'', ``Mary'', ``Fred'' and
``Susie'' have in common. We can tentatively call these expressions \emph{proper
names} (\emph{names}) and note that they refer \emph{directly} to the
individuals that they symbolize (which can include any kind of ``person, place
or thing''). For all intents and purposes, we can think of them as symbolizing
\emph{the objects themselves} as opposed to some mental or linguistic
intermediates. Names are a special case of \emph{singular terms} which pick out
individuals for \emph{reference}: they are special because they refer to
\textbf{exactly one} item without ambuguity. We will talk more about other
singular terms when we discuss quantifiers and variables. Names will be
symbolized in FOL using lower case latin letters ($a,b,c,\dots$).

Second, there is the rest of the sentence.

\begin{argument}
\item[ ] $[\hspace{1cm}]_x$ is reading.
\end{argument}

We will refer to such ``incomplete sentences'' as \emph{predicates} and to the
``gaps'' (here labelled with an $x$) as \emph{slots}. Note that predicates
can be made by introducing slots into multiple locations of sentences, depending
on how many are appropriate to the predicate word (i.e. \emph{verb},
\emph{adjective} or \emph{common noun} in question. Consider this one which
could have a gap in any of three locations.

\begin{argument}
\item[ ] $[\hspace{1cm}]_x$ gave a book to Mary.
\item[ ] John gave $[\hspace{1cm}]_x$ to Mary.
\item[ ] John gave a book to $[\hspace{1cm}]_x$.  
\end{argument}

We will call a verb like ``give'' a \emph{three-place predicate}, and
generally refer to a predicate with $n$ gaps as an \emph{n-place predicate}.
When $n>1$ will will sometimes refer to these predicates as \emph{n-place
relations}. Predicates and relations will be symbolized in FOL using capital
latin letters ($A,B,C,\dots$).

\subsection{Introducing FOL Quantifiers}

Notice that one of the singular terms used in the last example is not a name:
``a book''. This is where the quantifiers come in. Supposing that the gaps
in our predicates can themselves be labeled with symbols (\emph{variables},
symbolized by lower-case latin letters $\dots,x,y,z$), then we can tentatively
express ``John gave a book to Mary'' using \emph{quantifiers} to \emph{bind}
these variables. Let $G$ symbolize \emph{give}, $j$ and $m$ symbolize ``John''
and ``Mary'', respectively, $B$ symbolize ``is a book'' and $x$ stand for a
variable in both $G$ and $B$ that is also \emph{bound} by the quantifier
\eexists.

\begin{argument}
\item[ ] $\eexists x(B(x)\eand G(j,x,m))$
\end{argument}

Pronounced in ordinary English, this FOL sentence says ``there is an x such that
x is a book and John gives x to Mary'' which is more or less what we understand
the original English sentence to mean.

We will discuss binding and variables in more detail as we formalize the logic.
For now, note that there are two distinct \emph{quantifiers} in first-order
logic.

\begin{definition}[FOL Quantifiers (Informal)]
  \label{def:fol-quantifiers-informal}
  First-order logic has two distinct quantifiers:
  \begin{enumerate}
  \item \textbf{Existential Quantifier (\eexists x):} expresses the English
    paraphrase ``for at least one x''$\dots$. English expressions used to
    indicate this kind of quantification include ``there is/are'' and ``some''. 
  \item \textbf{Universal Quantifer (\eforall x):} expresses the English
    paraphrase ``for all x''$\dots$. English expressions used to indicate this
    kind of quantification include ``every'', ``any'', ``all'', and ``each''.
    \footnote{Beware that these terms are not used with perfect consistently
      in ordinary English, especially when subjunctive contexts are involved.}
    %% examples of how "any" varies by context (p241)
  \end{enumerate}
  In first-order logic, every quantifier must be followed by a variable
  indicating which of the predicate gaps it \emph{binds} in the interpretation
  of a sentence. Assigning variables to quantifier and gaps is the main
  mechanism for disabiguating complex English sentences within the FOL
  framework (that is not already available in TFL). Often it is more art than
  science.
\end{definition}

%% justify the conventions for quantifier sentences. (p271-2)
In a sense, when we talk about quantifiers, our subject is really quantified
sentences: formulas with a subject and more than one predicate. We will build
the sentences that use these quantifiers in a consistent and conventional way,
from $(Ac \eif Bc)$ for universals and from $(Aa \eand Ba)$ for existentials.
The form of these sentences is important to how quantifiers work in first-order
logic. Adding first-order quantifiers to the language in this way makes it
possible to model a variety of the inferences explored in classical logic by
ancient and medievel scholars.

\begin{definition}[Square of Opposition in FOL]
  Just these two quantifiers and negation are enough to express all four corners
  of Aristotle's \emph{square of opposition}. \\
\twoitems{
  \begin{argument}
    \item[ ] Every A is B.
    \item[ ] $\eforall x(A(x)\eif B(x))$
    \item[ ] 
    \item[ ] 
    \item[ ] Not every A is B.
    \item[ ] $\enot\eforall x(A(x)\eif B(x))$
  \end{argument}
}{
  \begin{argument}
    \item[ ] No A is B.
    \item[ ] $\enot\eexists x(A(x)\eand B(x))$
    \item[ ] 
    \item[ ] 
    \item[ ] Some A is B.
    \item[ ] $\eexists x(A(x)\eand B(x))$
  \end{argument}  
}
\end{definition}

Similarly, by folding in negation, we can represent additional English
quantifiers such as \emph{only}.

\begin{argument}
\item[ ] Only As are Bs.
\item[ ] $\enot\eexists x(B(x)\eand\enot A(x))$
\end{argument}

%% TODO: examples from English (pp 243-247)

\subsection{Domains of Individuals}

One question that arises immediately when dealing with a quantifier like
\eforall is ``every what?'' To answer this question we make reference to a
\emph{domain}.

\begin{definition}[Domain (Informal)]
  \label{def:fol-domain-informal}
  In first-order logic, a \emph{domain} is the collection of objects that we
  are talking about.
  \begin{enumerate}
  \item in general, \emph{singular terms} refer to objects in the domain
  \item every \emph{proper name} must pick out \textbf{exactly} one object in
    the domain;
  \item each domain item (\emph{individual}) can be picked out by any number of
    proper names (including zero).
  \end{enumerate}
\end{definition}

As we will see when we begin working through the semantics and natural deduction
systems for FOL, the choice of domain is a crucial determinant of the behavior
of a first-order logic. Ordinary English is typically understood as
\emph{fixing} the domain of its quantifying expressions, when we say ``Bobby''
we are typically clear about which individuals are relevant to our conversation,
especially the one that we mean by ``Bobby''. For us, choice of domain is
important, but more art than science.

% TODO: the art of domain selection (pp285-286)
We will typically define our domains in one of two ways: \emph{extensionally},
where we list all of the individuals; and \emph{descriptively}, where we use
ordinary language (or, later, a metalogical formula) to explain the contents of
the domain.

%%% empty domains
%      AxPx entails ExPx for non-empty domains (p330)
%      Ex(Px|-Px) is a theorem but not for empty domains (p330)
%      "free logic" (p331)

Among the effects of the domain on FOL involves the absence of referential
material. There are two prominent cases: when a singular term is used, but
it does not refer to any objects in the domain; when a predicate is used, but
it does not apply to any objects in the domain. An expression of the first kind
is \emph{Santa Claus}; an expression of the second is \emph{even prime greater
than 3}. Both of these phenomena can result in unexplected results. To
amelieroate this we will prohibit names that do not refer and we will give
predicates that do not apply to any objects (\emph{empty predicates}) a special
treatment.

\begin{definition}[Empty Predicates]
  When \ebeta is an empty predicate, any sentence
  $\eforall x(\ebeta(x)\eif\dots)$ is (vacuously) true.
\end{definition}

%% binary vs unary quantifers (p240, pp251-252)
%Additionally, when we restict the domain to just people, we can simplify
%$\eforall x(P(x)\eif M(x))$ (\emph{Every person is mortal.}) to just
%$\eforall xM(x)$ (\emph{Everyone is mortal.}) But quite often, our domains
%will be too rich to support this, including not just people but also numerous
%places and artefacts as well.

\subsection{Quantifiers and Variable Choice}

When we symbolize sentences using names, the identity of the name symbol is
enough to indicate that two sentences are ``about'' the same object.

\begin{argument}
  \item[ ] $A(c)\eand B(c)$
\end{argument}

This can be true even when the sentences are separated by any number of others
that may or may not contain any examples of the name in question. The name,
quite simply, always refers to exactly the same thing just in virtue of being a
proper name. 

But when a quantifier is used, the specific sentences to which the quantifier
should be applied is a more delicate matter. The following sentences all express
slightly different things.

\twoitems{
  \begin{argument}
  \item[ ] If c is A, then c is B.
  \item[ ] For all x, if c is A, then x is B.
  \item[ ] For all x, if x is A, then x is B.
  \end{argument}
}{
  \begin{argument}
  \item[ ] $A(c)\eif B(c)$
  \item[ ] $\eforall x(A(c)\eif B(x))$
  \item[ ] $\eforall x(A(x)\eif B(x))$
  \end{argument}
}

\subsection{Multi-Place Predicates}

Not only do we need to carefully track the variables across predicates, we also
need to keep them straight within a single predicate. As we noted above, some
predicates (i.e. \emph{n-place predicates}) could have more than one gap to fill
when used in FOL sentences. The example we presented for the verb \emph{give} is
repeated here.

\begin{argument}
\item[ ] $[\hspace{1cm}]_x$ gave a book to Mary.
\item[ ] John gave $[\hspace{1cm}]_x$ to Mary.
\item[ ] John gave a book to $[\hspace{1cm}]_x$.  
\end{argument}

We refer to the number of slots in an n-place predicate as its \emph{arity}
\footnote{\emph{unary}, \emph{binary}, \emph{ternary}, \emph{4-ary}, etc.},
and to an n-place predicate as \emph{n-ary}. We can consider the relevant
three-place predicate with all its gaps to be something like:

\begin{center}
  $[\hspace{1cm}]_x$ gave $[\hspace{1cm}]_y$ to $[\hspace{1cm}]_z$.
\end{center}

symbolized in FOL as:

\begin{center}
  $G(x,y,z)$
\end{center}

The order of $x,y,z$ is meaningful and important, but arbitrary. We can assign
the labels in the order that they appear in the original sentence. Note
that this convention must be used with care, because the following English
sentence has the same meaning but presents its gaps in a different order.

\begin{center}
  $[\hspace{1cm}]_x$ gave $[\hspace{1cm}]_z$ $[\hspace{1cm}]_y$.
\end{center}

The \emph{role} that the terms play in the sentence is more important than their
English order when deciding how to assign terms to the FOL predicate.

\subsection{Multiple Quantifiers}

When more than one quantifier is used, the process of interpretation can get
somewhat complicated. There are three major reasons for this:

\begin{observation}[Complications for the Interpretation of Multiple
    Quantifiers]
  When we have more that one quantifier, interpretation is complicated by three
  main factors:\\
  \begin{enumerate}
  \item The order that quantifiers occur is relevant to the interpretation of
    FOL sentences.
  \item The specific choice of variables that the quantifiers apply to is
    relevant to the interpretation of FOL sentences.
  \item The TFL connectives, especially negation, support radically different
    interpretations, depending on how they are taken to interact with
    quantifiers in FOL sentences.
  \end{enumerate}
\end{observation}

Let's look at each of these in turn. First, note that changing the order of the
quantifiers can change the meaning of the sentence.

\begin{argument}
\item[ ] $\eforall x\eexists y L(x,y)$ = Everyone loves someone. (maybe each
  someone different)
\item[ ] $\exists y\forall x L(x,y)$ = There is someone that everyone loves.
  (maybe Sara Lee) 
\end{argument}

Second, the application of a quantifier can vary. We will use `(', `)', `[' and
`]' to indicate the so-called \emph{scope} of a quantifier--the portion of the
overall sentence to which the quantifier applies. Here we have four distinct
quantifiers ($\forall x,\exists z,\exists y,\exists z$) even though the variable
$z$ is used to symbolize two of them.

\begin{argument}
\item[ ] $\forall x[\exists z(D(z)\eand O(x,z)) \eif
  \exists y(\exists z(D(z)\eand O(y,z)) \eand F(x,y))]$
\end{argument}

Finally, the interaction of negation with quantifier scope is a critical to the
interpretation of quantification phenomena that go beyond the two basic
quantifiers. Recall, for example, how we symbolize the square of opposition. The
flexible use of negation and scope allows us to express the FOL sentences with
whichever of the quantifiers that we prefer.

\twoitems{
  \begin{argument}
    \item[ ] Every A is B.
    \item[ ] No A is not B.
    \item[ ] $\eforall x(A(x)\eif B(x))$
    \item[ ] $\enot\eexists x(A(x)\eand\enot B(x))$
    \item[ ] 
    \item[ ] 
    \item[ ] Not every A is B.
    \item[ ] Some A is not B.
    \item[ ] $\enot\eforall x(A(x)\eif B(x))$
    \item[ ] $\eexists x(A(x)\eand\enot B(x))$
  \end{argument}
}{
  \begin{argument}
    \item[ ] No A is B.
    \item[ ] Every A is not B.
    \item[ ] $\enot\eexists x(A(x)\eand B(x))$
    \item[ ] $\eforall x(A(x)\eif\enot B(x))$
    \item[ ] 
    \item[ ] 
    \item[ ] Some A is B.
    \item[ ] Not every A is not B.
    \item[ ] $\eexists x(A(x)\eand B(x))$
    \item[ ] $\enot\eforall x(A(x)\eand\enot B(x))$
  \end{argument}  
}

\subsection{Paraphrasing English for First-Order Logic}

As we did for TFL, we can use paraphrase as a stepping stone to complete
formalization in FOL. A few rules of thumb might help.

\begin{fitchrule}[Methods of Paraphrase]
  When paraphrasing English as a intermediary to first-order logic, some
  consistent methods can help guide the process:
  \begin{outline}[enumerate]
    \1 Move quantifiers forward to clarify their scope, using pronouns to mark
    their slot
    \2 Every good boy deserves fudge.
    \2 (Every good boy is such that) he deserves fudge.
    \1 Replace pronouns with variables early, to avoid ambiguity. This is
    especially helpful when more than one quantifier is used.
    \2 (Every good boy is such that) he deserves fudge.
    \2 (Every good boy x is such that) x deserves fudge.
    \1 Set up the subject and predicate as paraphrases before symbolizing.
    \2 (Every good boy x is such that) x deserves fudge.
    \2 (Every x is such that)(x is a good boy \eif x deserves fudge)
    \1 Formalize the quantifiers once their subject predicates are handled.
    \2 (Every x is such that)(x is a good boy \eif x deserves fudge)
    \2 (\eforall x)(x is a good boy \eif x deserves fudge)
    \1 Once all of the basic sentences are made of just an n-ary predicate and
    n names or variables, then we can convert them to FOL sentences, using the
    symbolization key.
    \2 (\eforall x)(x is a good boy \eif x deserves fudge)
    \2 (\eforall x)(Gx \eif Fx)
  \end{outline}
\end{fitchrule}

Consider this sequence of examples adapted from \cite{Smi20}. For all of the
examples, we will assume the following extensionally defined interpretation.

\begin{interpretation}[Interpretation for Examples (1)-(12)]
  The interpretation will have three names.
  \begin{argument}
  \item[m:] Mariah
  \item[n:] Naomi
  \item[o:] Oliver
  \end{argument}
  Its domain will be all people currently living. And its predicates are defined
  descriptively as follows.
  \begin{argument}
  \item[F:] \textcircled{1} is a man
  \item[G:] \textcircled{1} is a woman
  \item[L:] \textcircled{1} loves \textcircled{2}
  \item[M:] \textcircled{1} is married to \textcircled{2}
  \item[R:] \textcircled{1} is the child of \textcircled{2} and \textcircled{3}
  \end{argument}
\end{interpretation}

\subsubsection{Examples with One Quantifier}

\begin{argument}
\item[\textbf{(1)}] There is someone who loves Mariah, Naomi and Oliver.
\item[$\simeq$] (Somone is such that) they love Mariah, Naomi and Oliver.
\item[$\simeq$] (Some x is such that) x loves Mariah, Naomi and Oliver.
\item[$\simeq$] (\eexists x) x loves Mariah, Naomi and Oliver.
\item[$\simeq$] (\eexists x)( x loves Mariah \eand x loves Naomi
  \eand x loves Oliver)  
\item[$\simeq$] \eexists x(Lxm\eand Lxn\eand Lxo)
\end{argument}

\begin{argument}
\item[\textbf{(2)}] Whoever is loved by Oliver is loved by Mariah too.
\item[$\simeq$] (Everyone who is loved by Oliver is such that)
  Mariah loves them.
\item[$\simeq$] (Every x who is loved by Oliver is such that)
  Mariah loves x.
\item[$\simeq$] (Every x is such that)
  (if x is loved by Oliver then Mariah loves x)
\item[$\simeq$] (\eforall x)(Oliver loves x \eif Mariah loves x)
\item[$\simeq$] \eforall x(Lox\eif Lmx)
\end{argument}

\begin{argument}
\item[\textbf{(3)}] Every woman who loves Mariah is loved by Oliver.
\item[$\simeq$] (Every woman x who loves Mariah is such that)
  x is loved by Oliver.
\item[$\simeq$] (Every x is such that)
  (if x is a woman who loves Mariah, then Oliver loves x)
\item[$\simeq$] (\eforall x)
  ( x is a woman and x loves Mariah \eif Oliver loves x)  
\item[$\simeq$] \eforall x((Gx\eand Lxm)\eif Lox)
\end{argument}

\begin{argument}
\item[\textbf{(4)}] Mariah loves some woman who loves Oliver.
\item[$\simeq$] (Some woman x who loves Oliver is such that) Mariah loves x
\item[$\simeq$] (Some x is such that)
  (x is a woman who loves Oliver and Mariah loves x)
\item[$\simeq$] (\eexists x)
  (x is a woman and x loves Oliver \eand Mariah loves x)
\item[$\simeq$] (\eexists x)
  (x is a woman \eand x loves Oliver \eand Mariah loves x)
\item[$\simeq$] \eexists x(Gx \eand Lxo \eand Lmx)
\end{argument}

\begin{argument}
\item[\textbf{(5)}] Naomi and Oliver love any child of theirs.
\item[$\simeq$] (Every x who is a child of Naomi and Oliver is such that)
  (Naomi and Oliver love x)
\item[$\simeq$] (Every x is such that)
  (if x is a child of Naomi and Oliver then Naomi and Oliver love x)
\item[$\simeq$] (\eforall x)
  (x is a child of Naomi and Oliver \eif (Naomi loves x \eand Oliver loves x))
\item[$\simeq$] \eforall x(Rxno \eif (Lnx \eand Lox))
\end{argument}

\begin{argument}
\item[\textbf{(6)}] Oliver is Mariah's child.
\item[$\simeq$] (Some x is such that)(Oliver is Mariah and x's child)
\item[$\simeq$] \eexists xRomx
\end{argument}

\begin{argument}
\item[\textbf{(7)}] Mariah is married to no one.
\item[$\simeq$] (it is not the case that)(Mariah is married to someone)
\item[$\simeq$] (it is not the case that)
  ((there is some x such that) Mariah is married to x)
\item[$\simeq$] (\enot)((\eexists x) Mariah is married to x)
\item[$\simeq$] \enot\eexists xMmx
\end{argument}

\begin{argument}
\item[\textbf{(8)}] No one who loves Naomi loves Oliver.
\item[$\simeq$] (It is not the case that)(Someone who loves Naomi loves Oliver)
\item[$\simeq$] \enot(Some x such that)(x loves Naomi \eand x loves Oliver)
\item[$\simeq$] \enot(\eexists x)(x loves Naomi \eand x loves Oliver)
\item[$\simeq$] \enot\eexists x(Lxn \eand Lxo)
\end{argument}

%\item[$\simeq$] \eforall x(Lxn\eand Gn)
\begin{argument}
\item[\textbf{(9)}] Naomi is a woman everyone loves.
\item[$\simeq$] Everyone loves Naomi \eand Naomi is a woman
\item[$\simeq$] (Every x is such that)(x loves Naomi) \eand Naomi is a woman
\item[$\simeq$] (\eforall x)(x loves Naomi) \eand Naomi is a woman
\item[$\simeq$] \eforall xLxn\eand Gn
\end{argument}

\begin{argument}
\item[\textbf{(10)}] If everyone loves Naomi then Oliver does.
\item[$\simeq$] (Every x is such that)(If x loves Naomi then Oliver loves Naomi)
\item[$\simeq$] (\eforall x)(x loves Naomi \eif Oliver loves Naomi)
\item[$\simeq$] \eforall x(Lxn \eif Lon)
\end{argument}

\begin{argument}
\item[\textbf{(11)}] If someone loves Naomi then Oliver does.
\item[$\simeq$] Someone loves Naomi \eif Oliver loves Naomi
\item[$\simeq$] (Some x is such that)(x loves Naomi) \eif Oliver loves Naomi
\item[$\simeq$] (\eexists x)(x loves Naomi) \eif Oliver loves Naomi
\item[$\simeq$] \eexists xLxn \eif Lon
\end{argument}

\begin{argument}
\item[\textbf{(12)}] If anyone loves Naomi then Oliver does.
\item[$\simeq$] (Everyone is such that)
  (If they love Naomi then Oliver loves Naomi)
\item[$\simeq$] (\eforall x)(x loves Naomi \eif Oliver loves Naomi)
\item[$\simeq$] \eforall x(Lxn \eif Lon)
\end{argument}

\subsubsection{Examples with Multiple Quantifiers}

%% NB: (13)
%% (pp282-285)
\begin{argument}
\item[\textbf{(1)}] Every man loves someone or other.
\item[$\simeq$] (Every man is such that) he loves someone or other.
\item[$\simeq$] (Every x is such that)
  (x is man \eif x loves someone or other)
\item[$\simeq$] (\eforall x)
    (x is man \eif (Some y such that)(x loves y))
\item[$\simeq$] (\eforall x)(x is a man \eif (\eexists y)(x loves y) )
\item[$\simeq$] \eforall x(Fx\eif\eexists yLxy)
\end{argument}

\begin{argument}
\item[\textbf{(2)}] Some man loves a woman.
\item[$\simeq$] (Some man is such that)(he loves a woman)
\item[$\simeq$] (Some x is such that)(x is a man \eand x loves a woman)
\item[$\simeq$] (\eexists x)(x is a man \eand
  (some woman y is such that)(x loves y) )
\item[$\simeq$] \eexists x(x is a man \eand
  (\eexists y)(x is a woman \eand x loves y) )
\item[$\simeq$] \eexists x(Fx\eand\eexists y(Gx \eand Lxy))
\end{argument}

\begin{argument}
\item[\textbf{(3)}] Someone Oliver loves is loved by everyone Naomi loves.
\item[$\simeq$] (Someone Oliver loves is such that)
  (they are loved by everyone Naomi loves)
\item[$\simeq$] (Some x is such that)
  (Oliver loves x \eand x is loved by everyone Naomi loves)
\item[$\simeq$] (\eexists x)
  (Oliver loves x \eand (Everyone Naomi loves is such that) x is loved by them)
\item[$\simeq$] \eexists x(Oliver loves x \eand
  (Everyone y is such that)(Naomi loves y \eif x is loved by y) )
\item[$\simeq$] \eexists x(Oliver loves x \eand
  (\eforall y)(Naomi loves y \eif y loves x))
\item[$\simeq$] \eexists x(Lox \eand\eforall y(Lny\eif Lyx))
\end{argument}

%% +1 reading (p283)
\begin{argument}
\item[\textbf{(4)}] No woman loves every man.
\item[$\simeq$] (It is not the case that)
  (Some woman is such that)
  (She loves every man)
\item[$\simeq$] \enot(Some x is such that)(x is a woman \eand x loves every man)
\item[$\simeq$] \enot(\eexists x)(x is a woman \eand x loves every man)
\item[$\simeq$] \enot(\eexists x)(x is a woman \eand
  (every man is such that)(x loves him))
\item[$\simeq$] \enot\eexists x(x is a woman \eand
  (every y is such that)(y is a man \eif x loves y))
\item[$\simeq$] \enot\eexists x(x is a woman \eand\eforall y(y is a man
  \eif x loves y))
\item[$\simeq$] \enot\eexists x(Gx\eand\eforall y(Fy\eif Lxy))
\end{argument}

%% +2 readings (p283)
\begin{argument}
\item[\textbf{(5)}] No woman loves any man.
\item[$\simeq$] (It is not the case that)
  (Some woman is such that)
  (She loves some man)
\item[$\simeq$] \enot(Some x is such that)(x is a woman \eand x loves some man)
\item[$\simeq$] \enot(\eexists x)(x is a woman \eand
  (some man is such that)(x loves him) )
\item[$\simeq$] \enot\eexists x(x is a woman \eand
  (some y is such that)(y is a man \eand x loves y) )
\item[$\simeq$] \enot\eexists x(x is a woman \eand
  (\eexists y)(y is a man \eand x loves y) )
\item[$\simeq$] \enot\eexists x(Gx \eand\eexists y(Fy \eand Lxy))
\end{argument}

\begin{argument}
\item[\textbf{(6)}] Anyone who is married loves someone they aren't married to.
\item[$\simeq$] (Everyone who is married is such that)
  (they love someone they aren't married to)
\item[$\simeq$] (Every x is such that)
  (x is married \eif x loves someone x isn't married to)
\item[$\simeq$] (\eforall x)
  (x is married \eif (some y is such that)
  (x loves y \eand x isn't married to y))
\item[$\simeq$] \eforall x(x is married \eif (\eexists y)
  (x loves y \eand x isn't married to y))
\item[$\simeq$] \eforall x(x is married \eif\eexists y(Lxy\eand\enot Mxy))
\end{argument}

\begin{argument}
\item[\textbf{(7)}] A married man only loves women.
\item[$\simeq$] (Every married man is such that)(he only loves women)
\item[$\simeq$] (\eforall x)(x is a married man \eif x only loves women)
\item[$\simeq$] (\eforall x)((x is a man \eand x is married) \eif
  x only loves women)
\item[$\simeq$] \eforall x((x is a man \eand x is married) \eif
  (Every y is such that(x loves y \eif y is a woman)))
\item[$\simeq$] \eforall x( (x is a man \eand x is married) \eif
  (\eforall y)(x loves y \eif y is a woman) )
\item[$\simeq$] \eforall x( (x is a man \eand
  (some z is such that)(x is married to z) ) \eif
  (\eforall y)(x loves y \eif y is a woman) )
\item[$\simeq$] \eforall x( (x is a man \eand
  (\eexists z)(x is married to z) ) \eif
  (\eforall y)(x loves y \eif y is a woman) )
\item[$\simeq$] \eforall x((Fx \eand\eexists zMxz)\eif\eforall y(Lxy\eif y Gy))
\end{argument}

\begin{argument}
\item[\textbf{(8)}] No man loves anyone who loves Naomi.
\item[$\simeq$] (it is not the case that)(some man is such that)
  (he loves anyone who loves Naomi)
\item[$\simeq$] \enot(some x is such that)
  (x is man \eand x loves anyone who loves Naomi)
\item[$\simeq$] \enot(\eexists x)
  (x is man \eand (every y is such that)(y loves Naomi \eif x loves y))
\item[$\simeq$] \enot\eexists x(x is man \eand
  \eforall y(y loves Naomi \eif x loves y))
\item[$\simeq$] \enot\eexists x(Fx\eand\eforall y(Lyn\eif Lxy))
\end{argument}

\subsection{Basic Inference Steps for First-Order Logic}

Based on the inference patterns we see in arguments using them, we can
tentatively describe four new valid inference steps to add to the system
we developed for truth-functional logic.

\begin{observation}[Valid Quantificational Inference Steps]
  There are four valid inference steps that we can take using the quantifiers
  of first-order logic:
  \begin{outline}[enumerate]
  \1 \textbf{Universal Instantiation:} From a universal sentence, we can
  infer any particular instance.
  \1 \textbf{Universal Generalization:} From a particular sentence proven
  true of \emph{an arbitrarily selected} individual, we can infer the
  universal generalization of that sentence.
  \1 \textbf{Existential Generalization:} From any particular sentence, we can
  infer the existential generalization of it.
  \1 \textbf{Existential Instantiation:} From an existential sentence, we can
  infer the existential instantiation of that sentence for a given name,
  provided:
  \2 the name is not already used elsewhere in the proof
  \2 the existential sentence does not contain the name
  \end{outline}
\end{observation}

\subsection{Identity of Individuals}

%% Metalogic ??
% equivalence relation (p362)
% equivalence class (p362)
% identity is the smallest equivalence relation (p363)
% Leibniz's Law (p364-5)

We still have more to say about singular terms. One open issue is the question
of how to keep them distinct. As it stands, we know that a name always refers
to the same individual in the domain (and that that individual is indeed in the
domain) but we cannot guess whether two given names refer to one individual or
two, both alternatives are permitted. Similarly, we have no way of expressing
the so-called \emph{count quantifiers} such as \emph{at least two}.

To tackle both of these problems, we will introduce a special two-place
predicate ($=$) that we will interpret as \emph{identity}. In other words,
when we write $a=b$, we will take that to mean that $a$ and $b$ both refer to
the same individual in the domain. So we can symbolize the simple identity
sentence ``Cicero is Tully'' with the FOL sentence $c=t$. We will call this type
of identity \emph{numerical identity}: the two names talk about
\textbf{the exact same (one) thing}.

We can use the negation of identity sentences to express concepts like
``everyone else'': $\eforall x(x\neq c\eif A(c,x))$, where $x\neq y$ means
$\enot(x=y)$. Likewise, when we want to symbolize ``only'', ``at least two''
and ``at most two'', respectively, we can use the FOL sentences:

%\item[ ] $(A(c)\eand\eforall y(y\neq c\eif\enot A(y)))$
\begin{argument}
\item[ ] $A(c)\eand\enot\eexists x(x=c\eand A(x))$
\item[ ] $\eexists x(A(x)\eand \eexists y(A(y)\eand x\neq y))$
\item[ ] $\enot\eexists x\eexists y\eexists z[((A(x)\eand A(y)\eand A(z))\eand((x\neq y\eand x\neq z)\eand y\neq z)]$
\end{argument}

When we want to express concepts like \emph{exactly 2}, we need to treat it
as \emph{at least 2 \textbf{and} at most 2}, symbolizing it as a conjuntion of
the two.

\begin{argument}
\item[ ] $\eexists x\eexists y[(A(x)\eand A(y)\eand x\neq y]\eand\enot\eexists x\eexists y\eexists z[((A(x)\eand A(y)\eand A(z))\eand((x\neq y\eand x\neq z)\eand y\neq z)]$
\end{argument}

\subsection{Definite Descriptions}

Finally, there is a type of expression whose interpretation seems to be
sometimes best treated as a singular term and sometimes as a quantifier.
Consider the following example.

\begin{argument}
\item[ ] The dog is waiting to come in.
\end{argument}

There are two important steps in the interpretation of this sentence that center
on the use of the expression ``the dog'' here. First, we are using it like a
singular term to express a proposition about a particular dog that we have in
our domain. Second, we are using quantificational elements to clarify which of
the domain objects we have in mind: namely, the one and only dog in the domain.

Famously, Bertrand Russell argued that we should merge these two aspects of such
expressions (callled \emph{definite descriptions}) by treating ``the'' as a
type of quantifier. We will look at the other appoach first.

Suppose we just stipulate that ``the dog'' works like a special kind of name,
picking out its referent \emph{directly} and operating as a singular term
whenever it occurs in a sentence. Then we might be able to symbolize the
sentence about our waiting dog as:

\begin{argument}
\item[ ] $W(\iota x(D(x)))$
\end{argument}

Here the expression $\iota x(D(x))$ behaves as a singular term, filling the
gap in the predicate $W$. For Russell's approach, we can just treat \emph{the}
as an \emph{existential quantifier} with a \emph{uniqueness condition}. We've
already seen both of these devices. We can use them together to express our
example:

\begin{argument}
\item[ ] $\eexists x((D(x)\eand\enot\eexists y(x\neq y\eand D(x)))\eand W(x))$
\end{argument}

There are problems with both approaches. We will adopt Russell's, since the
$\iota$ operator introduces \emph{truth gaps}, which are beyond the scope of
this introduction.

There are three traditional ways to formulate Russell's formula for definite
descriptions.

\begin{argument}
\item[\textbf{(R1)}] \eexists x(Ax\eand\eforall y(Ay\eif x=y) \eand Bx)
\item[\textbf{(R2)}] \eexists x\eforall y((Ay\eiff x=y)\eand Bx)
\item[\textbf{(R3)}] ((\eexists xAx\eand\eforall x\eforall y((Ax\eand Ay)\eif x=y))\eand\eforall x(Ax\eif Bx))
\end{argument}

As we will see, they are all equivalent.

